\begin{abstract}
Video diffusion transformers generate high-quality videos, but inference is
typically instance-agnostic: the same frozen denoising process is used regardless
of the motion and appearance statistics visible in the conditioning frames.
Test-time adaptation offers a natural way to specialize generation to each test
video, but conventional adaptation mechanisms are mismatched to modern video
diffusion transformers (DiTs). Full-model tuning exposes billions of parameters
to a single short clip, while LoRA-style adapters introduce thousands to millions
of free parameters and can overfit within a few gradient steps.
We introduce \method, Adaptive Shared Timestep Embedding Efficient Residuals,
an ultra-lightweight test-time adaptation method for video diffusion
transformers. \method learns a compact residual in the timestep embedding
pathway and shares it across all transformer blocks. Each block's frozen
adaptive layer-normalization projection maps this shared residual into
block-specific modulation vectors, yielding a structured weight-tying mechanism
that adapts denoising behavior with only a tiny number of trainable parameters
and no architectural changes. The residual is optimized only on observed
conditioning frames using the standard denoising objective, applied during
generation, and discarded afterward.
On \backbone, a 13.6B-parameter video DiT, \method matches, but does not
improve on, no-TTA on standard-horizon Panda-70M and UCF-101 settings: at
999 videos, standard-horizon Panda FVD is unchanged within noise
($154.7\to153.4$); the long-horizon Panda
1000-video evaluation shows a small distributional regression that we report
honestly as an open case rather than a success. These results suggest that
pretrained modulation pathways provide a strong substrate for efficient
per-video adaptation, while leaving long-horizon coherence as a clear next
target.
\end{abstract}
